\documentclass[tikz,border=6pt]{standalone}
% xystyle.tex — XY 스테이지 다이어그램 공통 스타일 (% xystyle.tex — XY 스테이지 다이어그램 공통 스타일 (% xystyle.tex — XY 스테이지 다이어그램 공통 스타일 (\input{xystyle})
%   흰 바탕, 직각 상자, 회색 + 남색 1색 강조, 수식은 LaTeX (Computer Modern), 한글은 kotex
\usepackage{kotex}
\usepackage{amsmath,amssymb}
\usetikzlibrary{arrows.meta,positioning,calc,fit,backgrounds}
\definecolor{navy}{HTML}{1F3A5F}
\definecolor{navybg}{HTML}{E9EEF5}
\definecolor{ink}{HTML}{333333}
\definecolor{hwbg}{HTML}{F3F3F3}
\definecolor{mute}{HTML}{777777}
\tikzset{
  >={Stealth[length=2.0mm,width=1.6mm]},
  every node/.style={text=ink},
  blk/.style={draw=ink, line width=0.6pt, fill=white, align=center, inner sep=4pt,
              minimum height=1.0cm, font=\small},
  key/.style={blk, draw=navy, fill=navybg, line width=0.8pt},
  hw/.style={blk, fill=hwbg},
  opt/.style={blk, dashed, draw=mute},
  sig/.style={->, line width=0.6pt, draw=ink},
  optsig/.style={->, line width=0.6pt, dashed, draw=mute},
  lab/.style={font=\footnotesize, align=center, inner sep=1.5pt},
  note/.style={font=\footnotesize, align=left, text=mute},
  grp/.style={draw=mute, dashed, line width=0.5pt, inner sep=7pt},
  ttl/.style={font=\footnotesize\bfseries, text=navy, anchor=south west, inner sep=1pt},
}
\newcommand{\tp}[1]{\texttt{\footnotesize #1}}
)
%   흰 바탕, 직각 상자, 회색 + 남색 1색 강조, 수식은 LaTeX (Computer Modern), 한글은 kotex
\usepackage{kotex}
\usepackage{amsmath,amssymb}
\usetikzlibrary{arrows.meta,positioning,calc,fit,backgrounds}
\definecolor{navy}{HTML}{1F3A5F}
\definecolor{navybg}{HTML}{E9EEF5}
\definecolor{ink}{HTML}{333333}
\definecolor{hwbg}{HTML}{F3F3F3}
\definecolor{mute}{HTML}{777777}
\tikzset{
  >={Stealth[length=2.0mm,width=1.6mm]},
  every node/.style={text=ink},
  blk/.style={draw=ink, line width=0.6pt, fill=white, align=center, inner sep=4pt,
              minimum height=1.0cm, font=\small},
  key/.style={blk, draw=navy, fill=navybg, line width=0.8pt},
  hw/.style={blk, fill=hwbg},
  opt/.style={blk, dashed, draw=mute},
  sig/.style={->, line width=0.6pt, draw=ink},
  optsig/.style={->, line width=0.6pt, dashed, draw=mute},
  lab/.style={font=\footnotesize, align=center, inner sep=1.5pt},
  note/.style={font=\footnotesize, align=left, text=mute},
  grp/.style={draw=mute, dashed, line width=0.5pt, inner sep=7pt},
  ttl/.style={font=\footnotesize\bfseries, text=navy, anchor=south west, inner sep=1pt},
}
\newcommand{\tp}[1]{\texttt{\footnotesize #1}}
)
%   흰 바탕, 직각 상자, 회색 + 남색 1색 강조, 수식은 LaTeX (Computer Modern), 한글은 kotex
\usepackage{kotex}
\usepackage{amsmath,amssymb}
\usetikzlibrary{arrows.meta,positioning,calc,fit,backgrounds}
\definecolor{navy}{HTML}{1F3A5F}
\definecolor{navybg}{HTML}{E9EEF5}
\definecolor{ink}{HTML}{333333}
\definecolor{hwbg}{HTML}{F3F3F3}
\definecolor{mute}{HTML}{777777}
\tikzset{
  >={Stealth[length=2.0mm,width=1.6mm]},
  every node/.style={text=ink},
  blk/.style={draw=ink, line width=0.6pt, fill=white, align=center, inner sep=4pt,
              minimum height=1.0cm, font=\small},
  key/.style={blk, draw=navy, fill=navybg, line width=0.8pt},
  hw/.style={blk, fill=hwbg},
  opt/.style={blk, dashed, draw=mute},
  sig/.style={->, line width=0.6pt, draw=ink},
  optsig/.style={->, line width=0.6pt, dashed, draw=mute},
  lab/.style={font=\footnotesize, align=center, inner sep=1.5pt},
  note/.style={font=\footnotesize, align=left, text=mute},
  grp/.style={draw=mute, dashed, line width=0.5pt, inner sep=7pt},
  ttl/.style={font=\footnotesize\bfseries, text=navy, anchor=south west, inner sep=1pt},
}
\newcommand{\tp}[1]{\texttt{\footnotesize #1}}

% diagram_tags : ZeroKey 태그 배치 세 안 (바디 좌표, 위에서 본 그림, 축척 1 m = 12 cm)
\begin{document}
\begin{tikzpicture}
\newcommand{\robot}[1]{%
  \begin{scope}[shift={(#1,0)}, scale=12]
    \fill[hwbg] (-0.21,-0.1715) rectangle (0.21,0.1715);
    \draw[ink, line width=0.6pt] (-0.21,-0.1715) rectangle (0.21,0.1715);
    \fill[mute!35] (-0.21,0.135) rectangle (0.21,0.1715);
    \fill[mute!35] (-0.21,-0.1715) rectangle (0.21,-0.135);
    \draw[navy, dashed, line width=0.6pt] (-0.0508,-0.0508) rectangle (0.0508,0.0508);
    \draw[->, ink, line width=0.6pt] (0.10,-0.11) -- (0.19,-0.11) node[right, font=\footnotesize, inner sep=1pt] {$x_B$};
    \draw[->, ink, line width=0.6pt] (0.10,-0.11) -- (0.10,-0.04) node[above, font=\footnotesize, inner sep=1pt] {$y_B$};
    \fill[red!70!black] (0,0) circle (0.006);
    \node[font=\footnotesize, anchor=north, inner sep=2pt] at (0,-0.006) {$P$};
  \end{scope}}
\newcommand{\ztag}[4]{% x y 이름 위치
  \fill[ink] ($(#1,#2)+(-0.11,-0.11)$) rectangle ($(#1,#2)+(0.11,0.11)$);
  \node[font=\footnotesize, #4, inner sep=2pt] at (#1,#2) {#3};}
% (a) 현재
\robot{0}
\ztag{-0.96}{0.72}{$T_1$}{above}
\draw[sig, draw=navy] (-0.96,0.72) -- node[lab, below left, text=navy, pos=0.5]{$d$} (0,0);
\node[blk, draw=none, fill=none, anchor=north, font=\footnotesize] at (0,-2.25) {\textbf{(a) 현재: 로봇에 1 개}\\[3pt]
  $\hat P = \hat p_1 + R(\hat\psi)(d + \Delta s)$\\[2pt]
  $\sigma_e^2 \approx \sigma_p^2 + \ell_u^2\big(\sigma_\psi^2 + \tfrac{(\dot b_\psi T)^2}{3}\big)$\\[2pt]
  예측 1.25 mm,\ 시뮬레이션 1.19 mm};
% (b) 로봇 2 개, 펜을 가운데 두고 대칭
\robot{6.4}
\ztag{5.44}{0.72}{$T_1$}{above}
\ztag{7.36}{-0.72}{$T_2$}{below}
\draw[navy, line width=0.5pt, dashed] (5.44,0.72) -- (7.36,-0.72);
\node[font=\footnotesize, text=navy, anchor=west] at (7.0,0.05) {$b = 0.2$ m};
\node[blk, draw=none, fill=none, anchor=north, font=\footnotesize] at (6.4,-2.25) {\textbf{(b) 로봇에 2 개 (펜을 가운데에)}\\[3pt]
  $\hat P = \tfrac12(\hat p_1 + \hat p_2) + R(\hat\psi)\,\Delta s$\\[2pt]
  $\sigma_e^2 \approx \tfrac12\sigma_p^2 + \lvert\Delta s\rvert^2\sigma_\psi^2$\\[2pt]
  센서 항 0.71 mm,\ 근사 시뮬레이션 0.84 mm};
% (c) 로봇 1 개 + 펜 캐리지 1 개
\robot{12.8}
\ztag{11.84}{0.72}{$T_1$}{above}
\ztag{12.8}{0.36}{$T_3$}{left=3pt}
\draw[->, navy, line width=0.5pt] (13.2,0.36) -- (13.55,0.36);
\draw[->, navy, line width=0.5pt] (13.2,0.36) -- (13.2,0.70);
\node[font=\footnotesize, text=navy, anchor=west] at (13.6,0.36) {캐리지와 함께};
\node[blk, draw=none, fill=none, anchor=north, font=\footnotesize] at (12.8,-2.25) {\textbf{(c) 로봇 1 개 + 펜 캐리지 1 개}\\[3pt]
  $\hat P = \hat p_3 + R(\hat\psi)\,r_3$\ \ ($r_3$: 태그 $\to$ 펜 끝)\\[2pt]
  $\sigma_e^2 \approx \sigma_p^2 + \lvert r_3\rvert^2\sigma_\psi^2$\\[2pt] 센서 항 1.0 mm,\ 근사 시뮬레이션 1.07 mm\\[2pt]
  스테이지를 거치는 되먹임 (지연 50 ms)};
% 범례
\node[note, anchor=north west] at (-2.6,-5.0) {회색 띠: 궤도.\ \ 남색 점선 사각형: 스테이지 행정 $\pm 50.8$ mm (펜 $P$ 둘레).\ \ 검은 사각형: ZeroKey 태그.\ \
  $T_1 = r_{tag} = (-0.08,\ 0.06)$ m,\ $T_2 = (0.08,\ -0.06)$ m.\\
  센서 항은 식 (22) 의 확장 ($\sigma_p = 1$ mm, $\sigma_\psi = 0.3^{\circ}$, 두 태그 잡음 독립 가정). 근사 시뮬레이션: (b) ZeroKey 잡음 $\sigma_p/\sqrt2$ + IMU 영향 0, (c) IMU 영향 0.\\ 센서 외 오차 (요철·크리프·스테이지 동특성) 바닥 약 0.41 mm 가 더해짐. 두 태그가 같은 다중 경로 오차를 받으면 (b) 의 이득은 줄어듦};
\end{tikzpicture}
\end{document}
